% answers_short.tex —— 2010 年·普通高等学校招生考试·湖北卷·理科 · 简短答案速查
% 由 tools/make_short_answers.py 自 answers.tex 的 \gksolhead/\gkanswer 提取，
% 并与 2010/2010_数学真题_湖北理_答案要点.md 逐题比对后生成，勿手改（改源 MD 后重跑）。

\begin{gknotice}
\textbf{使用说明}：本节只给最终答案便于快速核对，完整推导见解析版。答案为编写者独立推导，\textbf{非官方标准答案}。
\end{gknotice}

\gkpart{选择题}
\begin{center}\begin{tabular}{ccccccccccc}
\toprule
题号 & 1 & 2 & 3 & 4 & 5 & 6 & 7 & 8 & 9 & 10 \\
\midrule
答案 & D & A & D & C & B & B & C & B & C & A \\
\bottomrule
\end{tabular}\end{center}

\gkpart{填空题}
\begin{itemize}[leftmargin=2.6em]
  \item[\textbf{11.}] $6$。
  \item[\textbf{12.}] $5$。
  \item[\textbf{13.}] $4$。
  \item[\textbf{14.}] $0.4$。
  \item[\textbf{15.}] $CD$（几何平均数），$DE$（调和平均数）。
\end{itemize}

\gkpart{解答题}
\begin{itemize}[leftmargin=2.6em]
  \item[\textbf{16.}] （1）$\pi$；（2）最大值 $\dfrac{\sqrt2}{2}$，$x$ 的集合为 $\left\{x\ \middle|\ x=k\pi-\dfrac\pi8,\ k\in\mathbf Z\right\}$。
  \item[\textbf{17.}] （1）$k=40$，$f(x)=6x+\dfrac{800}{3x+5}$（$0\leqslant x\leqslant10$）；（2）$x=5$ 时总费用最小，最小值为 $70$ 万元。
  \item[\textbf{18.}] （1）证明见解析，$\dfrac{AB}{AQ}=3$；（2）$\dfrac{\sqrt{15}}5$。
  \item[\textbf{19.}] （1）$y^2=4x$（$x>0$）；（2）存在，$m\in\left(3-2\sqrt2,\,3+2\sqrt2\right)$。
  \item[\textbf{20.}] （1）$a_n=(-1)^{n-1}\sqrt{1-\dfrac34\left(\dfrac23\right)^{n-1}}$，$b_n=\dfrac14\left(\dfrac23\right)^{n-1}$；（2）证明见解析。
  \item[\textbf{21.}] （1）$b=a-1$，$c=1-2a$；（2）$a\in\left[\dfrac12,+\infty\right)$；（3）证明见解析。
\end{itemize}
